\chapter{Học sinh hỏi về thi cử và chuyện ôn thi}
\markboth{Học sinh hỏi về thi cử và chuyện ôn thi}{Học sinh hỏi về thi cử và chuyện ôn thi}

\section{Nước Tới Cổ Mới Nhảy}
\begin{enumerate}[leftmargin=*]
\item \hsqa{Thầy ơi bài này có thi không?}{Có thể có, nhưng cái đáng thi nhất là trí nhớ em đó.}{Teacher, will this be on the exam?}{It might be, but what's really being tested is your memory.}
\item \hsqa{Nếu không thi thì em khỏi học được không?}{Được, rồi kiến thức cũng khỏi vào em luôn.}{If it's not on the exam, can I skip learning it?}{Yes. And the knowledge will skip you entirely.}
\item \hsqa{Thầy ơi ôn bao nhiêu là đủ?}{Đủ là lúc em bớt hỏi câu này trong hoảng loạn.}{Teacher, how much revision is enough?}{Enough is when you stop asking this question in a panic.}
\item \hsqa{Em học trúng tủ được không?}{Được, nếu đề thương em như người nhà.}{Can I just guess the exam topics to study?}{You can, if the exam loves you like family.}
\item \hsqa{Em chỉ học phần dễ thôi.}{Đề thường không ký cam kết hiền như vậy.}{I'm only studying the easy parts.}{Exams usually don't sign a contract to be that gentle.}
\item \hsqa{Thi xong em quên sạch có sao không?}{Sao chứ, não em đâu phải phòng trọ theo ngày.}{Does it matter if I forget everything after the exam?}{It does. Your brain isn't a motel room booked by the day.}
\item \hsqa{Em ôn từ đêm trước có kịp không?}{Kịp lo, còn kịp giỏi thì khó.}{Is studying from the night before enough time?}{Enough time to panic. Enough time to get good is unlikely.}
\item \hsqa{Thầy đoán đề giúp em đi.}{Thầy dạy học, không hành nghề bói bài thi.}{Please guess the exam questions for me.}{I teach. I do not practice exam fortune-telling.}
\item \hsqa{Em làm đề cũ là đủ chưa?}{Đủ để biết mình còn thiếu bao nhiêu.}{Is doing past papers enough?}{Enough to know how much you are still lacking.}
\item \hsqa{Thi thử thấp quá em nản.}{Thi thử mà nản thì thi thật em tính làm thơ à?}{My mock exam score is so low I'm discouraged.}{Discouraged by a mock? What's your plan for the real one, writing poetry?}
\item \hsqa{Thầy ơi em run quá.}{Run là bình thường, miễn tay em đừng run tới mức quên não.}{Teacher, I'm trembling so much.}{Trembling is normal. Just don't let your hands tremble so much they forget the brain.}
\item \hsqa{Em học mà vào phòng thi là trắng đầu.}{Vậy em ôn kiến thức ít, ôn bản lĩnh còn ít hơn.}{I study but my mind goes blank in the exam hall.}{So you revised your knowledge a little, and your courage even less.}
\item \hsqa{Sao càng gần thi em càng muốn ngủ.}{Cơ thể em đang phản đối kiểu ôn nước rút đó.}{Why is it the closer to the exam, the sleepier I get?}{Your body is protesting that cramming method.}
\item \hsqa{Em làm đề mà điểm cứ dậm chân.}{Vì em đang luyện đề, chưa chắc đã luyện lỗi.}{I do practice papers but my score is stagnant.}{Because you are practicing tests, but not necessarily practicing fixing mistakes.}
\item \hsqa{Em chỉ cần qua môn thôi.}{Qua kiểu này thì kiến thức cũng lết qua cùng em.}{I just need to pass the subject.}{Passing like this means your knowledge is just dragging along too.}
\item \hsqa{Em thi xong mới học được không?}{Học sau thi nghe giống bơi sau khi chìm.}{Can I study after the exam is over?}{Studying after the exam sounds like learning to swim after you drown.}
\item \hsqa{Em sợ không đủ thời gian làm bài.}{Vậy lúc ôn em tập làm đúng giờ đi, đừng tập than.}{I'm afraid I won't have enough time to finish.}{Then practice working under a timer while revising. Don't practice complaining.}
\item \hsqa{Thi trắc nghiệm chắc dễ hơn.}{Dễ với người biết, khó với người đoán bằng niềm tin.}{Multiple choice exams must be easier.}{Easier for those who know. Harder for those guessing by pure faith.}
\item \hsqa{Em học thuộc hết là ổn chứ thầy?}{Thuộc mà không hiểu thì đề đổi tí là em tắt.}{If I memorise everything, I'll be fine, right?}{Memorising without understanding means if the question changes slightly, you shut down.}
\item \hsqa{Thầy ơi em ôn mà không thấy vào.}{Vì em đang quét mắt, chưa chắc đã động não.}{Teacher, I'm revising but nothing is going in.}{Because your eyes are scanning, but your brain isn't necessarily engaged.}
\end{enumerate}

\section{Ôn Thi Kiểu Lướt Sóng Và Những Hy Vọng Mù}
\begin{enumerate}[leftmargin=*,resume]
\item \hsqa{Em nên ôn từ chương nào trước?}{Từ chỗ em yếu nhất, đừng từ chỗ em thích nhất.}{Which chapter should I revise first?}{From where you are weakest. Not from where you like best.}
\item \hsqa{Em có nên thức trắng ôn thi không?}{Có thể, nếu em muốn hôm sau thi bằng linh hồn.}{Should I pull an all-nighter to revise?}{You can, if you want your soul to take the exam tomorrow.}
\item \hsqa{Em học nhóm mà toàn nói chuyện.}{Nhóm em đang ôn thi hay ôn tình bạn?}{I study in a group but we just talk.}{Is your group revising for the exam or revising for friendship?}
\item \hsqa{Em xem tóm tắt thôi được không?}{Được, miễn điểm em cũng thích sống tóm tắt.}{Can I just look at the summary?}{Sure, as long as your marks also like living summarily.}
\item \hsqa{Em không biết ôn kiểu gì cho nhớ.}{Lặp lại, làm lại, sửa lại; chứ không phải cầu lại.}{I don't know how to revise to remember.}{Repeat it, redo it, fix it. Not ''pray for it.''}
\item \hsqa{Em còn ít thời gian quá.}{Ít còn hơn không, đừng phí luôn phần ít đó.}{I have so little time left.}{Little is better than none. Don't waste the little part too.}
\item \hsqa{Em làm đề sai nhiều nên ngại làm tiếp.}{Sai nhiều lúc ôn còn rẻ hơn sai ít lúc thi.}{I make many mistakes on practice papers so I don't want to continue.}{Making many mistakes during revision is cheaper than making a few during the exam.}
\item \hsqa{Thầy ơi em có số thi cử không?}{Có, số lần em chịu ngồi xuống ôn đó.}{Teacher, is there such a thing as exam destiny?}{Yes. It's the number of times you're willing to sit down and revise.}
\item \hsqa{Em hay học xong quên lúc đi ngủ.}{Vậy chắc hôm sau em phải thức dậy ôn lại.}{I often study then forget it when I sleep.}{Then you'd better wake up the next day and revise again.}
\item \hsqa{Em thấy bạn em chưa học mà bình tĩnh lắm.}{Có khi bạn em gan lớn, chứ chưa chắc não lớn.}{My friend hasn't studied but is very calm.}{Maybe your friend has big guts, but not necessarily a big brain.}
\item \hsqa{Em đọc nhiều mà không thuộc.}{Đọc nhiều khác với đọc có mục tiêu.}{I read a lot but can't memorise it.}{Reading a lot is different from reading with a goal.}
\item \hsqa{Thi mà trúng câu khó thì sao thầy?}{Thì em làm câu dễ trước, chứ đừng ngồi nhìn đời trôi.}{What if the exam has really hard questions?}{Then you do the easy ones first. Don't just sit and watch life pass by.}
\item \hsqa{Em nên ôn dài hay ôn ngắn?}{Ôn thật thì kiểu nào cũng được, ôn giả thì kiểu nào cũng hỏng.}{Should I revise in long sessions or short?}{If it's genuine, either works. If it's fake, both fail.}
\item \hsqa{Em sợ mở đề ra thấy lạ.}{Lạ là đề làm việc, không phải đề phản bội em.}{I'm afraid I'll open the exam and find it alien.}{Alien is the exam doing its job, not betraying you.}
\item \hsqa{Em làm sai một câu là mất tinh thần.}{Tinh thần em mong manh hơn đáp án chuẩn luôn.}{I get one question wrong and lose my morale.}{Your morale is more fragile than an answer key.}
\item \hsqa{Em nên học một môn cả ngày không?}{Nếu não em chịu nổi thì được, còn thường nó phản ứng bằng cách lơ mơ.}{Should I study one subject all day?}{If your brain can endure it, yes. Usually, it responds by zoning out.}
\item \hsqa{Em chỉ học phần giáo viên nhấn mạnh.}{Nhấn mạnh không có nghĩa phần còn lại tự biến mất.}{I only studied the parts the teacher emphasized.}{Emphasized does not mean the rest of the syllabus deletes itself.}
\item \hsqa{Thầy ơi em ôn mà thấy vô vọng.}{Vô vọng thật là lúc em ngừng ôn.}{Teacher, I feel hopeless while revising.}{True hopelessness happens when you stop revising.}
\item \hsqa{Em nên xem lời giải trước hay tự làm trước?}{Tự vật lộn trước, kẻo não em quen ngồi xem phim.}{Should I look at the solution first or try it myself?}{Wrestle with it yourself first, lest your brain gets used to just watching movies.}
\item \hsqa{Em thấy đề mẫu dễ hơn đề thầy cho.}{Đời cũng hay vậy, bản demo thường hiền hơn bản chính thức.}{I found the sample paper easier than your assigned one.}{Life is often like that. The demo is usually gentler than the official release.}
\end{enumerate}

\section{Tâm Tình Phòng Thi Và Sự Đánh Rơi Phong Độ}
\begin{enumerate}[leftmargin=*,resume]
\item \hsqa{Em hay nhầm công thức lúc thi.}{Vì em nhớ mặt nó chứ chưa cưới nó vào đầu.}{I often mix up formulas during the exam.}{Because you memorize its face without marrying it to your mind.}
\item \hsqa{Em cứ ôn là lo.}{Lo vừa thôi để học, lo quá thì chỉ đủ mệt.}{I feel anxious whenever I revise.}{Anxious enough to study is good. Too anxious just makes you tired.}
\item \hsqa{Em nên bỏ bớt chương nào?}{Bỏ thì dễ, nhưng điểm nó sẽ nhớ em rất kỹ.}{Which chapters should I skip?}{Skipping is easy, but the lost marks will remember you very well.}
\item \hsqa{Em làm được ở nhà mà thi lại trượt.}{Ở nhà em thi với hòa bình, ở phòng thi em thi với áp lực.}{I can do it at home but fail out in the exam.}{At home you test with peace. In the hall you test with pressure.}
\item \hsqa{Em hay mất bình tĩnh khi canh giờ.}{Vậy lúc luyện em phải tập thêm đồng hồ, không chỉ tập bài.}{I lose my cool when working to a timer.}{Then when you practice, you must train with a clock, too.}
\item \hsqa{Em có nên xem mẹo khoanh đáp án không?}{Có, sau khi em đã học đàng hoàng, không phải thay cho học.}{Should I look up multiple choice tricks?}{Yes. After you've studied properly, not as a replacement.}
\item \hsqa{Em học từ khóa là đủ chứ?}{Từ khóa không tự nối thành hiểu biết đâu em.}{Learning keywords is enough, right?}{Keywords don't automatically connect themselves into comprehension.}
\item \hsqa{Em bị tâm lý nên thi thấp.}{Tâm lý cũng là một môn em đang nợ.}{My test anxiety caused my low score.}{Psychology is another subject you currently owe marks to.}
\item \hsqa{Em ôn mấy hôm mà quên sạch bài cũ.}{Thế nên mới có chuyện ôn vòng, không ai ôn một phát mà thành thánh.}{I study for a few days and forget the old stuff.}{That's why cyclic revision exists. No one becomes a saint on the first pass.}
\item \hsqa{Em làm bài chậm hơn mọi người.}{Chậm mà chắc còn đỡ hơn nhanh mà lao xuống hố.}{I work slower than everyone else.}{Slow and steady is better than fast and diving into a ditch.}
\item \hsqa{Thầy ơi em có cửa không?}{Có, miễn em đừng tự đóng trước.}{Teacher, do I have a chance?}{Yes, provided you don't lock the door yourself.}
\item \hsqa{Em chỉ sợ phụ lòng ba mẹ.}{Vậy ôn vì trách nhiệm, đừng chỉ ôn vì sợ.}{I'm just afraid of disappointing my parents.}{Then revise out of responsibility, not just fear.}
\item \hsqa{Em tính học ngày mai thật đó.}{Ngày mai rất đông người chen vào rồi.}{I honestly plan to study tomorrow.}{Tomorrow is already very crowded with other people's plans.}
\item \hsqa{Em nên luyện bao nhiêu đề?}{Đủ để em biết lỗi mình lặp ở đâu, không phải đủ để khoe số.}{How many practice papers should I do?}{Enough to know where you repeat mistakes. Not enough just to brag about the number.}
\item \hsqa{Em ôn xong mà vẫn không yên tâm.}{Yên tâm vừa đủ mới tỉnh, yên tâm quá là toang.}{I finished revising but still don't feel secure.}{Feeling secure enough keeps you sharp. Too secure and you're doomed.}
\item \hsqa{Em có nên bỏ ngủ để học không?}{Đừng đổi trí nhớ lấy quầng thâm.}{Should I give up sleep to study?}{Don't trade your memory for dark circles.}
\item \hsqa{Em mở đề lên mà tim rơi.}{Tim rơi thì nhặt lại, đề không nhặt hộ đâu.}{I open the exam paper and my heart drops.}{If it drops, pick it up. The paper won't do it for you.}
\item \hsqa{Thi xong em thấy mình làm ngu quá.}{Biết ngu sau thi cũng hơn tưởng mình khôn mãi.}{After the exam, I felt I did so stupidly.}{Realising that after is still better than thinking you're clever forever.}
\item \hsqa{Em nên chữa đề kiểu nào?}{Kiểu không tự tha cho những lỗi em lặp hoài.}{How should I review my practice papers?}{By never forgiving yourself for the errors you repeatedly make.}
\item \hsqa{Em mong đề đừng khó.}{Đề cũng đang mong em đừng chủ quan.}{I hope the paper isn't hard.}{The paper is hoping you aren't complacent.}
\end{enumerate}

\section{Chạy Trời Không Khỏi Áp Lực Thi Cử}
\begin{enumerate}[leftmargin=*,resume]
\item \hsqa{Em chỉ cần may thôi thầy.}{May là gia vị, không phải món chính.}{I just need luck, teacher.}{Luck is a spice, not the main course.}
\item \hsqa{Em học mà cứ thấy thiếu.}{Thiếu thì lấp dần, chứ hoảng không bù được.}{I study but always feel I'm missing something.}{If it's missing, gradually fill it. Panic doesn't compensate.}
\item \hsqa{Em ôn nhiều quá nên rối.}{Vì em gom nhiều nhưng chưa sắp.}{I reviewed so much I got confused.}{Because you gathered a lot but haven't sorted it.}
\item \hsqa{Em nên xem lại lý thuyết hay lao đề?}{Thiếu gốc thì quay về lý thuyết, đừng chạy như gà.}{Should I revisit theory or jump into practice tests?}{If your roots are weak, go back to theory. Don't run aimlessly.}
\item \hsqa{Em thi xong quên hết là bình thường không?}{Bình thường, miễn trước khi thi em từng thật sự nhớ.}{Is it normal to forget everything after the exam?}{Normal, as long as you were actually remembering before.}
\item \hsqa{Em ghét cảm giác đếm ngược đến ngày thi.}{Ghét cũng không làm lịch chậm hơn.}{I hate the exam countdown feeling.}{Hating it doesn't make the calendar move any slower.}
\item \hsqa{Em cứ nghe chữ thi là mệt.}{Mệt vì chữ thi hay vì mấy tuần trước em sống hơi liều?}{I feel exhausted just hearing the word 'exam'.}{Exhausted by the word, or because you lived recklessly the past few weeks?}
\item \hsqa{Em ôn ít nhưng chất lượng có được không?}{Được, nếu thật sự chất lượng chứ không phải nghe cho sang.}{Is it okay to study a little but with high quality?}{Fine, if it's genuinely high quality and not just a fancy excuse.}
\item \hsqa{Thầy ơi em sợ làm không kịp phần cuối.}{Vậy luyện cách phân thời gian, đừng chỉ luyện cầu nguyện.}{Teacher, I'm afraid I won't finish the final part.}{Then practice time management, not just praying.}
\item \hsqa{Em cần ai đó ngồi canh em học.}{Tốt, nhưng về lâu dài em phải tự canh nổi mình.}{I need someone to sit and monitor me studying.}{Good. But long term, you must monitor yourself.}
\item \hsqa{Em thi thử thấp quá, chắc thi thật toang.}{Thi thử thấp là bản tin cảnh báo, không phải cáo phó.}{My mock score was too low. Real exam is probably doomed.}{A low mock score is a weather warning, not an obituary.}
\item \hsqa{Em nên nghỉ trước ngày thi không?}{Nghỉ để tỉnh thì được, nghỉ để hoảng thì thôi.}{Should I rest the day before the exam?}{Rest to keep sharp, yes. Rest to panic, no.}
\item \hsqa{Em chỉ muốn qua được kỳ này.}{Qua rồi nhớ mang theo chút kiến thức, đừng qua tay không.}{I just want to get through this semester.}{Get through, but bring some knowledge with you. Don't leave empty-handed.}
\item \hsqa{Em làm đề sai mà không dám chữa.}{Không chữa thì lỗi đó ở lại như hộ khẩu.}{I got it wrong but don't dare correct it.}{If you don't correct it, that mistake will take up permanent residence.}
\item \hsqa{Em nên học tới phút cuối không?}{Tới mức đầu còn dùng được, đừng học tới mức mắt vô hồn.}{Should I study until the very last minute?}{Until your head still functions. Don't study until your eyes go blank.}
\item \hsqa{Em sợ kết quả không như mong.}{Sợ thì làm phần của mình tử tế hơn.}{I'm afraid the outcome won't be what I want.}{If you're afraid, handle your part of the job more earnestly.}
\item \hsqa{Thi cử làm em muốn bỏ chạy.}{Chạy đâu thì kiến thức cũng đuổi theo bài thi thôi.}{Exams make me want to run away.}{Wherever you run, the knowledge will chase you right onto the paper.}
\item \hsqa{Em có thể cầu may song song với ôn không?}{Được, miễn đừng bỏ phần ôn để tập trung cầu.}{Can I pray for luck while studying?}{You can, as long as you don't skip the studying to focus on the praying.}
\item \hsqa{Em thấy ôn thi là cực hình.}{Cực hình lớn hơn là vào phòng thi với đầu rỗng.}{I see revising as torture.}{A greater torture is entering an exam hall with an empty head.}
\item \hsqa{Thầy chốt giúp em: ôn thi cần nhất gì?}{Bớt ảo tưởng, tăng lặp lại, giữ bình tĩnh.}{Teacher, conclude for me: what is most important for exam prep?}{Less illusion, more repetition, keep your composure.}
\end{enumerate}

\section{Linh Cảm Bài Thi Và Nỗi Sợ Hãi Cuối Cùng}
\begin{enumerate}[leftmargin=*,resume]
\item \hsqa{Em có nên tin vào linh cảm khi làm bài?}{Tin vừa thôi, đáp án không chấm theo tâm linh.}{Should I trust my intuition during the test?}{Trust it a little. The answer sheet doesn't grade on the spiritual spectrum.}
\item \hsqa{Em học xong mà nhìn đề vẫn sợ.}{Sợ là bình thường, miễn tay vẫn làm.}{I studied, but I'm still scared when I look at the paper.}{Fear is normal. Just make sure your hand keeps moving.}
\item \hsqa{Em cứ thiếu tự tin phút chót.}{Thiếu tự tin thì kiểm tra lại, đừng ngồi tự hù.}{I always lose my confidence at the last minute.}{If you lose confidence, double-check your work. Don't sit there scaring yourself.}
\item \hsqa{Em mong thầy cho đề dễ.}{Thầy cũng mong em ôn dễ thương hơn.}{I hope you set an easy paper.}{I also hope you revise in a more endearing manner.}
\item \hsqa{Em làm một bài sai là hoảng cả buổi.}{Một bài sai không đáng, hoảng mới đáng hại.}{I get one question wrong and I panic the whole session.}{One wrong question isn't a big deal. Panic is what really does the damage.}
\item \hsqa{Em ôn nhiều nhưng không biết đủ chưa.}{Lúc em làm đề đỡ run hơn và lỗi lặp ít hơn là có dấu hiệu rồi.}{I revised a lot but I don't know if it's enough.}{When you do a paper with less trembling and fewer repeated mistakes, that's the sign.}
\item \hsqa{Em nên ngủ hay nên học tiếp?}{Nếu não em đang đi lạc thì ngủ rồi mai học tiếp.}{Should I sleep or keep studying?}{If your brain is wandering, sleep and continue tomorrow.}
\item \hsqa{Em sợ đề ra chỗ em bỏ.}{Đó là cách đời dạy em đừng bỏ bừa.}{I'm afraid the exam will test what I skipped.}{That's life's way of teaching you not to skip casually.}
\item \hsqa{Em thích ôn với bạn cho vui.}{Vui vừa thôi, đề không cộng điểm tình đồng chí.}{I like revising with friends for fun.}{Have fun moderately. The exam doesn't award marks for comradeship.}
\item \hsqa{Em cứ thi là đau bụng.}{Cơ thể em phản ứng hơi chân thật với cách em ôn.}{My stomach aches every time I take an exam.}{Your body is reacting quite honestly to your revision strategy.}
\item \hsqa{Em ôn cả tháng mà vẫn hoang mang.}{Hoang mang không chứng minh em yếu, chỉ chứng minh em biết kỳ thi quan trọng.}{I revised for a whole month and am still bewildered.}{Bewildered doesn't prove you're weak. It just proves you know the exam matters.}
\item \hsqa{Em sợ mình phụ công dạy của thầy.}{Vậy dùng nỗi sợ đó để ôn tử tế chứ đừng để than.}{I'm afraid I'll let down your teaching efforts.}{Then use that fear to revise properly instead of whining.}
\item \hsqa{Em làm câu nào cũng thấy nghi nghi.}{Nghi vừa đủ tốt, nghi quá thì tự phá nhịp.}{Every question I do, I feel suspicious.}{Doubt just enough is good. Doubt too much and you break your own rhythm.}
\item \hsqa{Em có nên bỏ câu khó không?}{Bỏ tạm để quay lại, đừng ôm rồi chết chung.}{Should I skip the hard questions?}{Skip them temporarily to return to them. Don't hug them and die together.}
\item \hsqa{Em cứ học là nghĩ tới trượt.}{Nghĩ ít lại, làm nhiều lên.}{Whenever I study, I think about failing.}{Think less. Do more.}
\item \hsqa{Em có nên đọc lại hết vở không?}{Nếu biết chọn trọng tâm thì đọc, không thì em đang ngắm chữ.}{Should I re-read my whole notebook?}{If you can pick the key points, yes. Otherwise, you're just sightseeing words.}
\item \hsqa{Em không quen áp lực.}{Vậy kỳ thi này vừa là thi kiến thức, vừa là tập thần kinh.}{I'm not used to pressure.}{Then this exam is partly a test of knowledge, and partly a workout for your nerves.}
\item \hsqa{Em muốn một câu an ủi trước khi thi.}{An ủi đây: chuẩn bị tử tế vẫn là cửa an toàn nhất.}{I want a comforting phrase before the exam.}{Here's comfort: adequate preparation is always the safest option.}
\item \hsqa{Em học gì cuối cùng trước giờ thi?}{Học bình tĩnh trước, rồi học lại lỗi hay sai.}{What should I revise right before the exam?}{Revise being calm first, then revise the mistakes you often make.}
\item \hsqa{Thầy ơi kỳ thi dạy em điều gì?}{Dạy rằng não phải làm việc trước khi may mắn xuất hiện.}{Teacher, what does an exam teach me?}{It teaches that your brain has to work before luck can make an appearance.}
\end{enumerate}